\section{Kontrola dowodu i najczęstsze pułapki}
\subsection{Łańcuch zależności}
Każda strzałka w poniższym schemacie została w tekście uzasadniona:
\[
\begin{gathered}
\text{słowa zredukowane}\xrightarrow{\text{tw. \ref{tw:redukcja}, stw. \ref{st:grupa-slow}}}F_2
\xrightarrow{\text{tw. \ref{tw:cztery-kawalki}}}2F_2,\\
F_2\xrightarrow{\text{tw. \ref{tw:wolne-obroty}}}G\le\mathrm{SO}(3)
\xrightarrow{\text{lem. \ref{lem:wyjatkowy}, wybór }T}X,\\
X\xrightarrow{\text{tw. \ref{tw:paradoksalnosc-x}}}2X,
\qquad S^2\xrightarrow{\text{lem. \ref{lem:lancuch}, \ref{lem:obrot-rozdzielajacy}}}X,\\
S^2\xrightarrow{\text{tw. \ref{tw:pelna-sfera}}}2S^2,\\
\text{stożki (lem. \ref{lem:stozki}):}\qquad B^3_*\equiv2B^3_*,\\
B^3\xrightarrow{\text{lem. \ref{lem:wchloniecie-srodka}}}B^3_*
\equiv2B^3_*\equiv2B^3
\xrightarrow{\text{przesunięcia}}\text{dwie rozłączne kule w }\R^3.
\end{gathered}
\]
\subsection{Dziesięć miejsc wymagających szczególnej uwagi}
\begin{enumerate}
\item \textbf{Redukcja nie polega na przestawianiu liter.} Usuwamy tylko sąsiadujące odwrotności. Jednoznaczność wyniku została udowodniona, a nie przyjęta intuicyjnie.
\item \textbf{Słowo puste musi należeć do partycji.} Włącza je zbiór $P$ w konstrukcji \eqref{eq:4-3}.
\item \textbf{Dwa niekomutujące obroty nie wystarczają.} Rachunek modulo $3$ wyklucza wszystkie niepuste słowa zredukowane.
\item \textbf{Nie redukujemy ułamka $1/3$ modulo $3$.} Redukujemy całkowite macierze liczników $N_s$.
\item \textbf{Grupa wolna to nie działanie swobodne.} Po usunięciu punktów ze zbioru $D$ wszystkie pozostałe stabilizatory są równe $\{e\}$.
\item \textbf{Wybór reprezentantów nie jest automatyczny.} Wskazaliśmy użycie aksjomatu wyboru przed lematem~\ref{lem:wspolrzedne-orbitowe}.
\item \textbf{Przeliczalne wyjątki nie znikają same.} Wchłania je dwuczęściowy rozkład oparty na obrocie spoza $G$.
\item \textbf{Kawałek złożony z nieskończenie wielu poziomów to nadal jeden kawałek.} Na całej unii poziomów $L$ działa jeden obrót $R$.
\item \textbf{Środek kuli wymaga osobnego ruchu.} Obroty wokół $0$ nigdy go nie przemieszczą; użyliśmy obrotu wokół osi przesuniętej w \eqref{eq:8-1}.
\item \textbf{Formalne kopie nie są jeszcze dwiema kulami w przestrzeni.} Dopiero końcowe przesunięcia o $v_1,v_2$ realizują etykiety jako rzeczywiste rozłączne położenia.
\end{enumerate}
\subsection{Krótkie zadania kontrolne z rozwiązaniami}
\begin{enumerate}
\item \textbf{Do której części należy $a^5$?} Należy do $P$, więc do $A_2$, a nie do $A_1$. Po ruchu $a$ daje $a^6$ w pierwszej kopii. Źródłem $e$ w tej kopii jest $\bar a$.
\item \textbf{Dlaczego z $hgt=hku$ wynika $g=k$ i $t=u$?} Najpierw skracamy działanie $h$, otrzymując $gt=ku$. Punkty $t,u$ reprezentują tę samą orbitę, zatem są równe. Wtedy $(k^{-1}g)t=t$ i swoboda działania daje $k^{-1}g=e$. Bez swobody ostatni wniosek nie byłby uzasadniony.
\item \textbf{Gdzie w rachunku macierzowym wykrywa się redukcję?} Dokładnie w zerowych wpisach tabeli \eqref{eq:5-5}. Sąsiednia para $ss^{-1}$ daje czynnik $w_s^Tw_{s^{-1}}=0$. Brak takich par gwarantuje, że wszystkie czynniki skalarne w \eqref{eq:5-7} są niezerowe.
\item \textbf{Dlaczego nie wystarczy policzyć samej liczby punktów?} Sama bijekcja $B^3\to2B^3$ nie musi być realizowana skończonym rozkładem, w którym każdy kawałek przenosimy ruchem sztywnym. W twierdzeniu istotne są jednocześnie: skończona liczba kawałków, rozłączność i ograniczenie do izometrii.
\item \textbf{Czy możemy powtórzyć podwojenie?} Tak. Po pierwszym podwojeniu stosujemy ten sam rozkład do każdej kuli osobno i otrzymujemy cztery kopie, potem osiem itd. Każda ustalona liczba powtórzeń nadal daje skończony rozkład, bo mnożymy tylko skończone liczby kawałków. Nie twierdzimy, że graniczna procedura z nieskończenie wielu powtórzeń jest skończonym rozkładem w sensie definicji~\ref{def:rownorozklad}.
\end{enumerate}

